\editionPage[\pdfbookmark[0]{Défis d’approfondissement}{bm-defis-debut-parent}]{defis-debut}{Dix-huit défis}{Un défi par chapitre · Parcours facultatif}{Chercher avant de consulter les deux indications.}
\begin{blocEdition}{defis-debut}
Ces problèmes demandent davantage d’initiative. Ils utilisent les outils du livre ; les prolongements sont définis dans l’énoncé. Aucun défi n’ajoute une exigence officielle au brevet.\par\bigskip
\begin{savoirfaire}[Une recherche qui avance]Faire un schéma, tester un petit cas, comparer deux méthodes, conserver les essais et préciser ce qui reste à prouver. Ouvrir la première indication seulement après une recherche réelle ; revenir au problème avant de lire la seconde.\end{savoirfaire}
\begin{savoirfaire}[Après la solution]Identifier l’idée décisive et la refaire sur une feuille vierge. Expliquer pourquoi elle fonctionne compte autant que retrouver le résultat.\end{savoirfaire}
\end{blocEdition}
\editionPage{defi-ex-1}{D01 · Défi}{Une fraction entre deux autres}{Approfondissement · 30 à 60 minutes indicatives}
\begin{blocEdition}{defi-ex-1}
Trouver une fraction de dénominateur 60 strictement comprise entre $5/12$ et $7/15$. Combien y en a-t-il ? Puis trouver une fraction comprise entre ces deux nombres, mais de dénominateur 120.
\end{blocEdition}
\editionPage{defi-ex-2}{D02 · Défi}{Comparer sans grands calculs}{Approfondissement · 30 à 60 minutes indicatives}
\begin{blocEdition}{defi-ex-2}
Comparer $2^{30}$ et $10^9$ sans développer les deux nombres. On pourra utiliser $2^{10}=1024$. Puis dire si $2^{60}$ est plus grand que $10^{18}$.
\end{blocEdition}
\editionPage{defi-ex-3}{D03 · Défi}{Une recherche entière exhaustive}{Approfondissement · 30 à 60 minutes indicatives}
\begin{blocEdition}{defi-ex-3}
Trouver tous les entiers positifs n tels que n divise 180 et que n soit un carré parfait. Justifier que la liste est complète.
\end{blocEdition}
\editionPage{defi-ex-4}{D04 · Défi}{L’identité de Sophie Germain}{Approfondissement · 30 à 60 minutes indicatives}
\begin{blocEdition}{defi-ex-4}
Vérifier que pour tous nombres $a$ et $b$ : $a^4+4b^4=(a^2+2b^2)^2-(2ab)^2$. En déduire une factorisation de $n^4+4$, puis démontrer que $n^4+4$ n’est jamais un nombre premier lorsque $n$ est un entier supérieur ou égal à 2.
\end{blocEdition}
\editionPage{defi-ex-5}{D05 · Défi}{Deux solutions et un contexte}{Approfondissement · 30 à 60 minutes indicatives}
\begin{blocEdition}{defi-ex-5}
Un rectangle a pour côtés x et $x+3$ cm. Un carré a pour côté $x+1$ cm. Pour quelles valeurs positives de x leurs aires diffèrent-elles de 1 cm$^2$, sans préciser laquelle est la plus grande ?
\end{blocEdition}
\editionPage{defi-ex-6}{D06 · Défi}{Promotions dans quel ordre ?}{Approfondissement · 30 à 60 minutes indicatives}
\begin{blocEdition}{defi-ex-6}
Un article coûte P euros, avec $P>20$. Le magasin applique une réduction de 20 \% et un bon de 10 €. Comparer les deux ordres possibles et montrer que la différence de prix ne dépend pas de P.
\end{blocEdition}
\editionPage{defi-ex-7}{D07 · Défi}{Une rencontre de débits}{Approfondissement · 30 à 60 minutes indicatives}
\begin{blocEdition}{defi-ex-7}
Une cuve vide reçoit 18 L/min et perd simultanément 3 L/min par une fuite constante. Après 12 min, on coupe l’arrivée d’eau ; la fuite continue seule. Au bout de combien de temps depuis le début la cuve redevient-elle vide ?
\end{blocEdition}
\editionPage{defi-ex-8}{D08 · Défi}{Une fonction cachée derrière des essais}{Approfondissement · 30 à 60 minutes indicatives}
\begin{blocEdition}{defi-ex-8}
On sait seulement que $f(-1)=3$, $f(0)=2$, $f(1)=3$. Un élève conclut que $f(x)=x^2+2$. Montrer que cette conclusion n’est pas imposée en proposant une autre formule donnant les mêmes trois images mais une image différente pour 2.
\end{blocEdition}
\editionPage{defi-ex-9}{D09 · Défi}{Trois tarifs, un choix argumenté}{Approfondissement · 30 à 60 minutes indicatives}
\begin{blocEdition}{defi-ex-9}
Pour n séances entières, de 1 à 20, trois offres coûtent A = 8n, B = $30+5n$, C = $70+2n$, en euros. Déterminer l’offre la moins chère pour chaque nombre de séances ; indiquer les égalités utiles.
\end{blocEdition}
\editionPage{defi-ex-10}{D10 · Défi}{Une série à reconstruire}{Approfondissement · 30 à 60 minutes indicatives}
\begin{blocEdition}{defi-ex-10}
Proposer une série de huit nombres entiers dont le minimum est 2, le premier quartile 5, la médiane 8, le troisième quartile 12, le maximum 20 et la moyenne 9. Une telle série peut-elle avoir une moyenne égale à 15 ?
\end{blocEdition}
\editionPage{defi-ex-11}{D11 · Défi}{Un jeu de somme équitable ?}{Approfondissement · 30 à 60 minutes indicatives}
\begin{blocEdition}{defi-ex-11}
Deux dés équilibrés sont lancés indépendamment. Le joueur A gagne si la somme est 6 ou 8 ; B gagne si elle est 7. Sinon on recommence sans modifier les règles. Les deux joueurs ont-ils la même chance de gagner lors d’un lancer décisif ? Comparer en comptant les couples, sans nouveau formalisme.
\end{blocEdition}
\editionPage{defi-ex-12}{D12 · Défi}{Une corde dans un cercle}{Approfondissement · 30 à 60 minutes indicatives}
\begin{blocEdition}{defi-ex-12}
Un cercle a pour centre O et rayon 13 cm. Une corde AB est à 5 cm de O. H est le pied de la perpendiculaire issue de O sur AB. On admet que H est le milieu de la corde. Calculer AB et l’aire de OAB.\label{fig-defi-12}\figureGeo{defi-corde}

\end{blocEdition}
\editionPage{defi-ex-13}{D13 · Défi}{Les diagonales d’un trapèze}{Approfondissement · 30 à 60 minutes indicatives}
\begin{blocEdition}{defi-ex-13}
ABCD est un trapèze de bases [AB] et [CD], parallèles, avec AB = 6 cm et CD = 10 cm. Ses diagonales [AC] et [BD] se coupent en O. Calculer $OA/OC$. Si AC = 12 cm, calculer OA et OC.\label{fig-defi-13}\figureGeo{defi-trapeze}

\end{blocEdition}
\editionPage{defi-ex-14}{D14 · Défi}{Deux angles de visée}{Approfondissement · 30 à 60 minutes indicatives}
\begin{blocEdition}{defi-ex-14}
Deux points d’observation A et B sont à 10 m l’un de l’autre sur une même droite horizontale allant vers le pied d’une tour. B est le plus proche. Les angles d’élévation du sommet sont de 30 degrés en A et 45 degrés en B. Les visées se font au niveau du sol dans ce modèle. Calculer la hauteur au décimètre.\label{fig-defi-14}\figureGeo{defi-visees}

\end{blocEdition}
\editionPage{defi-ex-15}{D15 · Défi}{Une somme de deux vecteurs}{Approfondissement · 30 à 60 minutes indicatives}
\begin{blocEdition}{defi-ex-15}
ABC est un triangle. On construit le point D tel que $\vect{AD}=\vect{AB}+\vect{AC}$, puis le point E tel que $\vect{BE}=\vect{AC}$. Démontrer que D et E sont confondus. Quelle est la nature de ABDC ? Si I est le milieu de [BC], montrer que $\vect{AB}+\vect{AC}=2\vect{AI}$.\label{fig-defi-15}\figureGeo{defi-vecteurs}

\end{blocEdition}
\editionPage{defi-ex-16}{D16 · Défi}{Une section d’aire imposée}{Approfondissement · 30 à 60 minutes indicatives}
\begin{blocEdition}{defi-ex-16}
Une boule de rayon 10 cm est coupée par un plan. Le disque de section doit avoir pour aire $64\pi$ cm$^2$. À quelle distance du centre placer le plan ? Y a-t-il plusieurs positions parallèles possibles ?\label{fig-defi-16}\figureGeo{defi-section}

\end{blocEdition}
\editionPage{defi-ex-17}{D17 · Défi}{Un script de recherche et sa preuve}{Approfondissement · 30 à 60 minutes indicatives}
\begin{blocEdition}{defi-ex-17}
Un script initialise s à 0, n à 1, puis répète tant que s est inférieur à 100 : ajouter n à s ; ajouter 1 à n. Quelle valeur de n possède-t-il à la sortie ? Écrire un script par blocs ou en français structuré et justifier l’arrêt.
\end{blocEdition}
\editionPage{defi-ex-18}{D18 · Défi}{Le meilleur rectangle entier}{Approfondissement · 30 à 60 minutes indicatives}
\begin{blocEdition}{defi-ex-18}
Un rectangle a des côtés entiers positifs et un périmètre de 46 cm. Déterminer les dimensions donnant l’aire maximale et prouver qu’aucun autre rectangle entier ne fait mieux.
\end{blocEdition}
